%! Graphing Rational Functions — Unit 4, Lesson 4.5 — Homework (Answer Key)
\documentclass[10pt]{article}
\usepackage{algebra2-article}
\usepackage{algebra2-key}   % pulls in -boxes; do NOT also load -boxes
\usepackage{grading}
\gradingassignment{alg2-u04-l05-hw}
\gradingcourse{algebra_2}
\usepackage{algebra2-graded}
\newcommand{\TallMath}[1]{$\displaystyle #1\rule[-1.4em]{0pt}{3.2em}$}
% Standard coordinate grid + axes: {xmin}{xmax}{ymin}{ymax}
\newcommand{\lgrid}[4]{%
  \draw[step=1, linegray!40, very thin] (#1,#3) grid (#2,#4);
  \draw[->, semithick, charcoal] (#1,0)--(#2,0) node[below right, font=\scriptsize]{$x$};
  \draw[->, semithick, charcoal] (0,#3)--(0,#4) node[left, font=\scriptsize]{$y$};}

\begin{document}

\pageheader{Unit 4, Lesson 4.5}{Homework --- Answer Key}

\vspace{0.10in}

\begin{notesbox}{Practice}
Every problem starts the same way: \textbf{factor}. Restrictions come off the \emph{original}
denominator, and the cancel test sorts each one into a hole or a wall. Write each final answer in its box. Only what is inside a box is graded.

\begin{enumerate}[leftmargin=1.9em, itemsep=10pt]
  \item \textbf{Hole or wall?} Give the restrictions, then sort them.
    \begin{enumerate}[label=(\alph*), leftmargin=2.1em, itemsep=5pt]
      \item $\dfrac{x+3}{(x+3)(x-5)}$ \hfill restrictions \inblank{\blank{1.8cm}}{\gradedblank[id=q1a, type=short, points=1, answer={-3, 5}, check=set, display={$-3,5$}]{1.8cm-2pt}{18pt}} \; hole \inblank{\blank{1.2cm}}{\gradedblank[id=q1b, type=short, points=1, answer={x=-3}]{1.2cm-2pt}{18pt}} \;
            VA \inblank{\blank{2.0cm}}{\gradedblank[id=q1c, type=short, points=1, answer={x=5}]{2.0cm-2pt}{10pt}}
      \item $\dfrac{x-6}{(x-2)(x+1)}$ \hfill restrictions \inblank{\blank{1.8cm}}{\gradedblank[id=q1d, type=short, points=1, answer={2, -1}, check=set, display={$2,-1$}]{1.8cm-2pt}{18pt}} \; hole \inblank{\blank{1.2cm}}{\gradedblank[id=q1e, type=short, points=1, answer={none}, check=text]{1.2cm-2pt}{18pt}} \;
            VA \inblank{\blank{2.0cm}}{\gradedblank[id=q1f, type=short, points=1, answer={x=2, x=-1}, check=set, display={$x=2$, $x=-1$}]{2.0cm-2pt}{10pt}}
      \item $\dfrac{x^2-25}{x^2-5x}$ \hfill restrictions \inblank{\blank{1.8cm}}{\gradedblank[id=q1g, type=short, points=1, answer={0, 5}, check=set, display={$0,5$}]{1.8cm-2pt}{18pt}} \; hole \inblank{\blank{1.2cm}}{\gradedblank[id=q1h, type=short, points=1, answer={x=5}]{1.2cm-2pt}{18pt}} \;
            VA \inblank{\blank{2.0cm}}{\gradedblank[id=q1i, type=short, points=1, answer={x=0}]{2.0cm-2pt}{10pt}}
      \item $\dfrac{x^2+x-6}{x^2-4}$ \hfill restrictions \inblank{\blank{1.8cm}}{\gradedblank[id=q1j, type=short, points=1, answer={2, -2}, check=set, display={$2,-2$}]{1.8cm-2pt}{18pt}} \; hole \inblank{\blank{1.2cm}}{\gradedblank[id=q1k, type=short, points=1, answer={x=2}]{1.2cm-2pt}{18pt}} \;
            VA \inblank{\blank{2.0cm}}{\gradedblank[id=q1l, type=short, points=1, answer={x=-2}]{2.0cm-2pt}{10pt}}
      \item $\dfrac{7}{x^2-16}$ \hfill restrictions \inblank{\blank{1.8cm}}{\gradedblank[id=q1m, type=short, points=1, answer={4, -4}, check=set, display={$4,-4$}]{1.8cm-2pt}{18pt}} \; hole \inblank{\blank{1.2cm}}{\gradedblank[id=q1n, type=short, points=1, answer={none}, check=text]{1.2cm-2pt}{18pt}} \;
            VA \inblank{\blank{2.0cm}}{\gradedblank[id=q1o, type=short, points=1, answer={x=4, x=-4}, check=set, display={$x=4$, $x=-4$}]{2.0cm-2pt}{10pt}}
      \item $\dfrac{x^2-9}{x^2-6x+9}$ \hfill restrictions \inblank{\blank{1.8cm}}{\gradedblank[id=q1p, type=short, points=1, answer={3}, check=set]{1.8cm-2pt}{18pt}} \; hole \inblank{\blank{1.2cm}}{\gradedblank[id=q1q, type=short, points=1, answer={none}, check=text]{1.2cm-2pt}{18pt}} \;
            VA \inblank{\blank{2.0cm}}{\gradedblank[id=q1r, type=short, points=1, answer={x=3}]{2.0cm-2pt}{10pt}}
    \end{enumerate}
    \vspace{2pt}
    Part (f) is the odd one. The factor $x-3$ appears on top \emph{and} on the bottom, and there is
    still no hole. Why not? \rubricbox{\gradedblank[id=q1s, type=rubric, points=2, display={the denominator has \emph{two} copies of $x-3$ and the top has one, so one copy survives the cancelling --- leaving $\frac{x+3}{x-3}$, a wall}]{\linewidth}{10pt}}

  \item \textbf{Horizontal asymptotes.} Degrees only --- do not factor.
    \begin{enumerate}[label=(\alph*), leftmargin=2.1em, itemsep=4pt]
      \item $y=\dfrac{5x-1}{x^2+3}$: \; HA \inblank{\blank{1.6cm}}{\gradedblank[id=q2a, type=short, points=1, answer={y=0}]{1.6cm-2pt}{18pt}}
      \item $y=\dfrac{2x^2+x}{3x^2-5}$: \; HA \inblank{\blank{1.6cm}}{\gradedblank[id=q2b, type=short, points=1, answer={y=2/3}, display={$y=\frac23$}]{1.6cm-2pt}{18pt}}
      \item $y=\dfrac{x^3+1}{x^2-4}$: \; HA \inblank{\blank{1.6cm}}{\gradedblank[id=q2c, type=short, points=1, answer={none}, check=text]{1.6cm-2pt}{18pt}}
      \item $y=\dfrac{-4x}{x-9}$: \; HA \inblank{\blank{1.6cm}}{\gradedblank[id=q2d, type=short, points=1, answer={y=-4}]{1.6cm-2pt}{18pt}}
    \end{enumerate}

  \item \textbf{Two full builds.}
    \begin{keepwithprompt}
    \begin{enumerate}[label=(\alph*), leftmargin=2.1em, itemsep=8pt]
      \item $f(x)=\dfrac{x+2}{x^2-2x-8}$ \\[3pt]
            Factored \inblank{\blank{3.0cm}}{\gradedblank[id=q3a, type=short, points=1, answer={(x+2)/((x-4)(x+2))}, display={$(x+2)/\bigl((x-4)(x+2)\bigr)$}]{3.0cm-2pt}{10pt}} \; Restrictions \inblank{\blank{2.0cm}}{\gradedblank[id=q3b, type=short, points=1, answer={4, -2}, check=set, display={$4,-2$}]{2.0cm-2pt}{10pt}} \; Simplified \inblank{\blank{1.8cm}}{\gradedblank[id=q3c, type=short, points=2, answer={1/(x-4)}, display={$1/(x-4)$}]{1.8cm-2pt}{10pt}}
            \\[3pt]
            Hole \inblank{\blank{1.8cm}}{\gradedblank[id=q3d, type=short, points=2, answer={(-2,-1/6)}, check=text, display={$(-2,-1/6)$}]{1.8cm-2pt}{10pt}} \; VA \inblank{\blank{1.2cm}}{\gradedblank[id=q3e, type=short, points=1, answer={x=4}]{1.2cm-2pt}{10pt}} \; HA \inblank{\blank{1.2cm}}{\gradedblank[id=q3f, type=short, points=1, answer={y=0}]{1.2cm-2pt}{10pt}} \;
            $x$-int \inblank{\blank{1.6cm}}{\gradedblank[id=q3g, type=short, points=1, answer={none}, check=text]{1.6cm-2pt}{10pt}} \; $y$-int \inblank{\blank{1.6cm}}{\gradedblank[id=q3h, type=short, points=1, answer={(0,-1/4)}, check=text, display={$(0,-1/4)$}]{1.6cm-2pt}{10pt}} \\[3pt]
            Sign chart --- boundary points \inblank{\blank{1.6cm}}{\gradedblank[id=q3i, type=short, points=1, answer={4}, check=set]{1.6cm-2pt}{10pt}}. \; Test $x=0$: \inblank{\blank{1.0cm}}{\gradedblank[id=q3j, type=short, points=1, answer={-}, check=text]{1.0cm-2pt}{10pt}} \;
            test $x=5$: \inblank{\blank{1.0cm}}{\gradedblank[id=q3k, type=short, points=1, answer={+}, check=text]{1.0cm-2pt}{10pt}} \\[3pt]
            Domain \inblank{\blank{5.0cm}}{\gradedblank[id=q3l, type=short, points=1, answer={(-inf,-2)U(-2,4)U(4,inf)}, check=text, display={$(-\infty,-2)\cup(-2,4)\cup(4,\infty)$}]{5.0cm-2pt}{10pt}}
      \item $g(x)=\dfrac{x-3}{x^2-4}$ \\[3pt]
            Factored \inblank{\blank{3.0cm}}{\gradedblank[id=q3m, type=short, points=1, answer={(x-3)/((x-2)(x+2))}, display={$(x-3)/\bigl((x-2)(x+2)\bigr)$}]{3.0cm-2pt}{10pt}} \; Restrictions \inblank{\blank{2.0cm}}{\gradedblank[id=q3n, type=short, points=1, answer={2, -2}, check=set, display={$2,-2$}]{2.0cm-2pt}{10pt}} \; Holes \inblank{\blank{1.4cm}}{\gradedblank[id=q3o, type=short, points=1, answer={none}, check=text]{1.4cm-2pt}{10pt}}
            \\[3pt]
            VA \inblank{\blank{2.4cm}}{\gradedblank[id=q3p, type=short, points=1, answer={x=2, x=-2}, check=set, display={$x=2$, $x=-2$}]{2.4cm-2pt}{10pt}} \; HA \inblank{\blank{1.2cm}}{\gradedblank[id=q3q, type=short, points=1, answer={y=0}]{1.2cm-2pt}{10pt}} \; $x$-int \inblank{\blank{1.6cm}}{\gradedblank[id=q3r, type=short, points=1, answer={(3,0)}, check=text]{1.6cm-2pt}{10pt}} \;
            $y$-int \inblank{\blank{1.6cm}}{\gradedblank[id=q3s, type=short, points=1, answer={(0,3/4)}, check=text, display={$(0,3/4)$}]{1.6cm-2pt}{10pt}} \\[3pt]
            Sign chart --- boundary points \inblank{\blank{2.4cm}}{\gradedblank[id=q3t, type=short, points=1, answer={-2, 2, 3}, check=set, display={$-2$, $2$, $3$}]{2.4cm-2pt}{10pt}}.
            \begin{center}
            \renewcommand{\arraystretch}{1.4}
            \begin{tabularx}{0.86\linewidth}{|C{2.4cm}|C{1.5cm}|C{1.6cm}|X|}
            \hline
            \rowcolor{forestbg}
            \textbf{Interval} & \textbf{Test $x$} & \textbf{Sign} & \textbf{Above / below} \\ \hline
            $x<-2$   & $-3$  & \inblank{\hspace{1.2cm}}{\gradedblank[id=q3u, type=short, points=1, answer={-}, check=text]{1.2cm-2pt}{10pt}} & \inblank{\hspace{2.6cm}}{\gradedblank[id=q3v, type=short, points=1, answer={below}, check=text]{2.6cm-2pt}{10pt}} \\ \hline
            $-2<x<2$ & $0$   & \inblank{\hspace{1.2cm}}{\gradedblank[id=q3w, type=short, points=1, answer={+}, check=text]{1.2cm-2pt}{10pt}} & \inblank{\hspace{2.6cm}}{\gradedblank[id=q3x, type=short, points=1, answer={above}, check=text]{2.6cm-2pt}{10pt}} \\ \hline
            $2<x<3$  & $2.5$ & \inblank{\hspace{1.2cm}}{\gradedblank[id=q3y, type=short, points=1, answer={-}, check=text]{1.2cm-2pt}{10pt}} & \inblank{\hspace{2.6cm}}{\gradedblank[id=q3z, type=short, points=1, answer={below}, check=text]{2.6cm-2pt}{10pt}} \\ \hline
            $x>3$    & $4$   & \inblank{\hspace{1.2cm}}{\gradedblank[id=q3aa, type=short, points=1, answer={+}, check=text]{1.2cm-2pt}{10pt}} & \inblank{\hspace{2.6cm}}{\gradedblank[id=q3ab, type=short, points=1, answer={above}, check=text]{2.6cm-2pt}{10pt}} \\ \hline
            \end{tabularx}
            \end{center}
            \vspace{-4pt}
            Domain \inblank{\blank{5.0cm}}{\gradedblank[id=q3ac, type=short, points=1, answer={(-inf,-2)U(-2,2)U(2,inf)}, check=text, display={$(-\infty,-2)\cup(-2,2)\cup(2,\infty)$}]{5.0cm-2pt}{10pt}} \; This graph crosses its horizontal asymptote at \inblank{\blank{1.6cm}}{\gradedblank[id=q3ad, type=short, points=1, answer={(3,0)}, check=text]{1.6cm-2pt}{10pt}}.
    \end{enumerate}
    \end{keepwithprompt}

  \item \textbf{Which equation is this?} The dashed lines are the asymptotes, the solid dot is an
        intercept, and the open circle is a hole.
    \begin{center}
    \begin{tikzpicture}[scale=0.36]
      \lgrid{-7}{7}{-5}{7}
      \draw[navy, dashed, semithick] (1,-5)--(1,7);
      \draw[navy, dashed, semithick] (-7,1)--(7,1);
      \node[navy, font=\scriptsize, above right] at (1,6.2) {$x=1$};
      \node[navy, font=\scriptsize, below left] at (-3.4,1) {$y=1$};
      \begin{scope}
        \clip (-7,-5) rectangle (7,7);
        \draw[very thick, forest, domain=-7:0.62, samples=140, smooth] plot (\x,{1+2/((\x)-1)});
        \draw[very thick, forest, domain=1.33:7, samples=140, smooth] plot (\x,{1+2/((\x)-1)});
      \end{scope}
      \draw[forest, very thick, fill=white] (-3,0.5) circle (4.0pt);
      \fill[charcoal] (-1,0) circle (3.2pt) node[above right, font=\scriptsize]{$(-1,0)$};
      \fill[charcoal] (0,-1) circle (3.2pt) node[below right, font=\scriptsize]{$(0,-1)$};
    \end{tikzpicture}
    \end{center}
    \vspace{-4pt}
    \begin{center}
    \textbf{A.} $\dfrac{x+1}{x-1}$ \quad
    \textbf{B.} $\dfrac{(x+1)(x+3)}{(x-1)(x+3)}$ \quad
    \textbf{C.} $\dfrac{(x+1)(x-3)}{(x-1)(x-3)}$ \quad
    \textbf{D.} $\dfrac{x+3}{(x-1)(x+3)}$
    \end{center}
    \vspace{3pt}
    Answer \inblank{\blank{1.0cm}}{\gradedblank[id=q4a, type=short, points=1, answer={B}, check=text]{1.0cm-2pt}{10pt}}. \; The hole is at \inblank{\blank{1.6cm}}{\gradedblank[id=q4b, type=short, points=1, answer={(-3,1/2)}, check=text, display={$(-3,1/2)$}]{1.6cm-2pt}{10pt}}, so the domain is \inblank{\blank{4.2cm}}{\gradedblank[id=q4c, type=short, points=1, answer={(-inf,-3)U(-3,1)U(1,inf)}, check=text, display={$(-\infty,-3)\cup(-3,1)\cup(1,\infty)$}]{4.2cm-2pt}{10pt}}.
    \\[3pt]
    Finally, rewrite the simplified form as a transformed parent (Lesson 4.4):
    $\dfrac{x+1}{x-1}=\dfrac{(x-1)+\inblank{\rule{0.8cm}{0.4pt}}{\gradedblank[id=q4d, type=short, points=1, answer={2}, check=numeric]{0.8cm-2pt}{8pt}}}{x-1}=\inblank{\rule{2.4cm}{0.4pt}}{\gradedblank[id=q4e, type=short, points=1, answer={1+2/(x-1)}, display={$1+\frac{2}{x-1}$}]{2.4cm-2pt}{18pt}}$
    \\[3pt]
    {\small \textbf{A} has no hole; \textbf{C} puts it at $+3$; \textbf{D} gives HA $y=0$.}

  \item \textbf{Error analysis.} A student analyzed $g(x)=\dfrac{x^2-x-6}{x^2-9}$ and wrote:
    \begin{center}
    \fbox{\parbox{0.84\linewidth}{\centering \vspace{2pt}
    ``It simplifies to $\dfrac{x+2}{x+3}$, so the domain is all reals except $x=-3$, \\
    and the graph is exactly the graph of $\dfrac{x+2}{x+3}$.''\vspace{2pt}}}
    \end{center}
    \begin{enumerate}[label=(\alph*), leftmargin=2.1em, itemsep=5pt]
      \item The simplification is \inblank{\blank{1.4cm}}{\gradedblank[id=q5a, type=short, points=1, answer={correct}, check=text]{1.4cm-2pt}{10pt}} (correct / incorrect).
      \item What did the student leave out of the domain, and why? \rubricbox{\gradedblank[id=q5b, type=rubric, points=2, display={$x=3$ --- the \emph{original} denominator $x^2-9$ is zero there, so $3$ was never an allowed input}]{\linewidth}{10pt}}
      \item The hole is at \inblank{\blank{1.8cm}}{\gradedblank[id=q5c, type=short, points=1, answer={(3,5/6)}, check=text, display={$(3,5/6)$}]{1.8cm-2pt}{10pt}}, so the correct domain is \inblank{\blank{4.4cm}}{\gradedblank[id=q5d, type=short, points=1, answer={(-inf,-3)U(-3,3)U(3,inf)}, check=text, display={$(-\infty,-3)\cup(-3,3)\cup(3,\infty)$}]{4.4cm-2pt}{10pt}}.
      \item State the rule from Lesson 4.1 that the student broke, in one sentence. \rubricbox{\gradedblank[id=q5e, type=rubric, points=2, display={cancelling never restores an input --- restrictions are read off the original denominator, and the cancelled one leaves a hole behind}]{\linewidth}{10pt}}
    \end{enumerate}

  \item \textbf{Cleaning the smokestack.} A factory must remove pollution from its exhaust. Removing
        $p$ percent of the pollutant costs
        \[ C(p)=\frac{25p}{100-p} \quad \text{thousand dollars,} \qquad 0\le p<100. \]
    \begin{enumerate}[label=(\alph*), leftmargin=2.1em, itemsep=5pt]
      \item $C(50)=\inblank{\blank{1.4cm}}{\gradedblank[id=q6a, type=short, points=1, answer={25}, check=numeric]{1.4cm-2pt}{10pt}}$, \; $C(80)=\inblank{\blank{1.4cm}}{\gradedblank[id=q6b, type=short, points=1, answer={100}, check=numeric]{1.4cm-2pt}{10pt}}$, \; $C(90)=\inblank{\blank{1.4cm}}{\gradedblank[id=q6c, type=short, points=1, answer={225}, check=numeric]{1.4cm-2pt}{10pt}}$, \;
            $C(95)=\inblank{\blank{1.4cm}}{\gradedblank[id=q6d, type=short, points=1, answer={475}, check=numeric]{1.4cm-2pt}{10pt}}$, \; $C(99)=\inblank{\blank{1.4cm}}{\gradedblank[id=q6e, type=short, points=1, answer={2475}, check=numeric]{1.4cm-2pt}{10pt}}$ (thousand dollars).
      \item Run the build. Restriction \inblank{\blank{1.6cm}}{\gradedblank[id=q6f, type=short, points=1, answer={p!=100}, display={$p\neq100$}]{1.6cm-2pt}{10pt}}; vertical asymptote \inblank{\blank{1.6cm}}{\gradedblank[id=q6g, type=short, points=1, answer={p=100}]{1.6cm-2pt}{10pt}};
            degrees \inblank{\blank{0.6cm}}{\gradedblank[id=q6h, type=short, points=1, answer={1}, check=numeric]{0.6cm-2pt}{10pt}} and \inblank{\blank{0.6cm}}{\gradedblank[id=q6i, type=short, points=1, answer={1}, check=numeric]{0.6cm-2pt}{10pt}}, so the horizontal asymptote is \inblank{\blank{1.6cm}}{\gradedblank[id=q6j, type=short, points=1, answer={y=-25}]{1.6cm-2pt}{10pt}}.
      \item Which of those two asymptotes means something about this factory, and what does it say?
            And why is the other one meaningless here? \gradedlines{2}{\gradedblank[id=q6k, type=rubric, points=3, display={the wall at $p=100$: as the target approaches $100\%$ the cost grows without bound, so total removal cannot be bought at any price. $y=-25$ describes $p$ running far past $100$, at a negative cost --- but $p$ only lives between $0$ and $100$}]{\linewidth}{\gradedlinesheight}}
      \item The $x$-intercept is \inblank{\blank{1.4cm}}{\gradedblank[id=q6l, type=short, points=1, answer={(0,0)}, check=text]{1.4cm-2pt}{10pt}}. In plain words, what does that say? \rubricbox{\gradedblank[id=q6m, type=rubric, points=1, display={removing none of it costs nothing}]{\linewidth}{10pt}}
      \item Going from $90\%$ to $95\%$ costs an extra \inblank{\blank{1.8cm}}{\gradedblank[id=q6n, type=short, points=1, answer={250}, check=numeric]{1.8cm-2pt}{10pt}} thousand dollars. Going from
            $95\%$ to $99\%$ costs an extra \inblank{\blank{1.8cm}}{\gradedblank[id=q6o, type=short, points=1, answer={2000}, check=numeric]{1.8cm-2pt}{10pt}} thousand. What is the graph doing near
            $p=100$, and what does that mean for a law requiring $100\%$ removal? \gradedlines{2}{\gradedblank[id=q6p, type=rubric, points=2, display={it climbs steeply along the vertical asymptote --- each additional percent costs more than the one before, and a $100\%$ requirement is a requirement no finite budget can meet}]{\linewidth}{\gradedlinesheight}}
    \end{enumerate}
\end{enumerate}
\end{notesbox}

\begin{extensionbox}
\textbf{Part 1 --- Design one.} Write the equation of a rational function with vertical asymptotes at
$x=-1$ and $x=4$, a horizontal asymptote $y=0$, and \emph{no} $x$-intercept at all.
\begin{center}
$y=$ \ans{$\dfrac{1}{(x+1)(x-4)}$}
\end{center}
\begin{enumerate}[label=(\alph*), leftmargin=2.1em, itemsep=5pt]
  \item Which requirement forced the denominator? \ans{the two vertical asymptotes --- and neither factor may cancel}
  \item Which two requirements, together, forced the numerator to be a nonzero \emph{constant}?
        \ansline{``no $x$-intercept'' means the numerator is never zero, and HA $y=0$ means its degree is below $2$ --- together that leaves a nonzero constant}
  \item Is your answer the only one? What could someone else have written instead? \ans{no --- any nonzero constant on top works, e.g.\ $\frac{5}{(x+1)(x-4)}$ or $\frac{-2}{x^2-3x-4}$}
\end{enumerate}

\vspace{4pt}
\textbf{Part 2 --- One question, carefully.} Can a rational function have a hole \emph{and} a vertical
asymptote at the same $x$-value? Answer, then test your answer against $y=\dfrac{x-3}{(x-3)^2}$ ---
which has $x-3$ on the top and on the bottom.
\ansline{No. The cancel test sends each restriction to exactly one outcome: the factor either cancels
completely (hole) or leaves a copy in the denominator (wall). In $\frac{x-3}{(x-3)^2}$ one copy cancels
and one remains, leaving $\frac{1}{x-3}$ --- a wall at $x=3$ and no hole. A shared factor is not
enough; it has to cancel \emph{all the way}.}
\end{extensionbox}

\begin{spiralbox}
\textbf{Coming next --- Lesson 4.6: Solving Rational Equations.} All unit long you have computed
excluded values and then \emph{kept} them off the graph. Tomorrow they come back as answers: clearing
the denominators hands you numbers that solve the cleared equation but not the original one --- exactly
the excluded values you have been listing since 4.1. You already know how to find them.
\end{spiralbox}

\begin{teachernote}
\textbf{1(f)}: two copies of $x-3$ downstairs, one upstairs --- the restriction survives as a wall and
there is no hole. \textbf{3(a)} has \emph{no} $x$-intercept; \textbf{3(b)} crosses its HA at $(3,0)$.
\textbf{Problem 5} is the unit's spine: the algebra is right and the conclusion still wrong. Accept any
phrasing of 4.1's \emph{cancelling never restores an input}, and flag anyone missing it --- it is
tomorrow's extraneous solution. \textbf{Problem 6} reverses the activity's model: here the
\emph{vertical} asymptote is the one that means something.
\end{teachernote}

\end{document}
